\section{Introduction}
\label{sec:intro}

Periodic elastic media can be designed to control wave propagation through their dispersion. Early calculations established band structures for periodic composites; experiments and theory later showed that two-dimensional solids can support absolute acoustic gaps \cite{SigalasEconomou1992,kushwaha1993,vasseur2001}. Lattice analyses and topology-optimization studies have since used geometry to tailor pass bands, stop bands, and wave transmission \cite{MartinssonMovchan2003,SigmundJensen2003,Dalklint2022}. Recent methods target gaps at specified frequencies and use multiple materials to control gap formation, while scaling studies track how optimized gaps shift with structural size \cite{He2026,Nakayama2025,Quinteros2026}. Reviews cover phononic-crystal mechanisms, engineered band gaps, and applications \cite{Vasileiadis2021,Oudich2023,MaWangYuan2026}. A gap along selected directions does not guarantee a complete gap; the latter requires suppression throughout the Brillouin zone.

When a wavelength approaches the material's microstructural scale, a local Cauchy law may not capture size-dependent stiffness or dispersion. Higher-order theories add kinematic measures and their work-conjugate stresses. Their development draws on couple-stress and microstructure theories, second-gradient elasticity, and variational treatments of higher-order boundary work \cite{toupin1962,MindlinTiersten1962,mindlin1964,Mindlin1965,mindlineshel1968,Germain1973}. Later studies examined strain-gradient descriptions of nonuniform deformation and small-scale beam experiments \cite{Aifantis1999,YangChongLamTong2002,LamYangChongWangTong2003}. Reviews and lattice-based derivations describe other formulations of gradient elasticity \cite{askesaifantis2011,polyzosfotiadis2012}. We use the local Mindlin Form-II theory; nonlocal integral and micromorphic theories are related approaches, but they are distinct formulations.

A dynamic gradient model needs a kinetic-energy law as well as gradient stiffness. With classical inertia alone, the shear branch examined here has phase speed proportional to wavenumber at high wavenumbers. Independent gradient micro-inertia gives this branch a finite high-wavenumber limit. That result follows from the constitutive and kinetic-energy laws used here; other higher-order continua need not have the same limit. One-dimensional studies have examined dispersion in gradient-elastic beams and periodic media, including dipolar-gradient formulations \cite{papargyribeskou2009,liweizhou2016,HosseiniSladekSladekZhang2024}. A two-dimensional phononic crystal has also been analyzed with a relaxed micromorphic model that includes free and gradient micro-inertia \cite{MadeoColletMiniaciBillonOuisseNeff2018}. Together, these studies show that the continuum and inertia laws affect the interpretation of short-wavelength bands, although their constitutive models differ from the local Mindlin Form-II model used here.

Size-dependent phononic-crystal studies use a range of constitutive models. Transfer-matrix analyses show that imperfect interfaces can alter one-dimensional band gaps \cite{zhengwei2009}. Other one-dimensional work considers dipolar-gradient elasticity, gradient thermoelasticity, and flexoelectric strain-gradient coupling \cite{liweizhou2016,li2023anchorA,li2024anchorB,HosseiniSladekSladekZhang2024}. In two dimensions, researchers have applied nonlocal strain-gradient theory to nanoscale phononic crystals and studied mixed and defect modes in a phononic-crystal slab \cite{JinHuHu2022,JinHuHu2023}. Other model-specific comparisons include thermoelastic cylindrical crystals and functionally graded lattices \cite{hosseinizhang2021,mishra2026anchorC}. Anisotropy and orientation have also been studied in conventional phononic crystals: work on orthotropic fillers and anisotropic plates reports shifts in band locations and widths as material axes or inclusions rotate relative to the lattice \cite{zhanwei2010,YaoWuHouXin2011}. Recent studies examine direction-dependent gaps in anisotropic nanocrystal assemblies and the effects of anisotropic constituents in optimized phononic crystals \cite{Kim2023,Liu2024anisotropic}. For flexural plate waves, anisotropic equal-frequency contours can align group velocity with a lens interface over a broad frequency range \cite{Danawe2023}. Prior work has therefore established orientation-dependent band tuning. Here we examine how an oriented characteristic-length tensor affects dispersion and energy-flow direction in a local two-dimensional strain-gradient model with gradient micro-inertia.

The finite-element approximation must also accommodate the strain-gradient weak form. Because it contains second derivatives of displacement, the trial space must be $H^2$-conforming or use an equivalent mixed formulation. The bicubic Bogner--Fox--Schmit (BFS) rectangle provides $C^1$ continuity across elements on structured quadrilateral meshes \cite{bfs1965}. Earlier studies developed displacement-based, mixed, and implicit-gradient finite elements \cite{ShuKingFleck1999,AmanatidouAravas2002,ZervosPapanicolopulosVardoulakis2009,Beheshti2017}. Later work treats second-gradient boundary-value problems, higher-order Hermite mappings on regular and distorted domains, mixed elements with patch recovery or shared degrees of freedom, and mixed crack formulations \cite{AndreausDellIsolaGiorgioPlacidiLekszyckiRizzi2016,BacciocchiFantuzzi2023,BacciocchiFantuzzi2025,ChoiLeeSim2023,Papanicolopulos2024,ChirkovNazarenkoAltenbach2024}. For a conforming BFS Bloch element, the same phase transformation must be applied to displacement values and all derivative degrees of freedom, including the mixed derivative.

The studies cited above do not combine, in one two-dimensional Bloch--Floquet finite-element framework, a rotated ellipsoidal second-moment characteristic-length tensor, local Mindlin Form-II strain-gradient elasticity with independent gradient micro-inertia, and a conforming $C^1$ BFS discretization with phase transformations on all relevant degrees of freedom. This study addresses that combination and builds on established work on anisotropy, orientation effects, and size-dependent phononic crystals. It makes four contributions:
\begin{itemize}
  \item We formulate a Bloch--Floquet framework for two-dimensional anisotropic strain-gradient elasticity. The rotated ellipsoidal averaging domain defines the length tensor, with aspect ratio $\mathrm{AR}=l_1/l_2$ and orientation $\theta$ as separate design variables.
  \item Gradient micro-inertia is examined analytically and numerically. For the branch studied, $\ell_{\mathrm{i}}=0$ gives an unbounded high-wavenumber phase speed, whereas $\ell_{\mathrm{i}}>0$ yields the finite limit $v_{T,\infty}=\sqrt{\mu/\rho}\,l_{\mathrm{eff}}/(\sqrt{10}\,\ell_{\mathrm{i}})$ derived in \ref{app:asymptotics}.
  \item A conforming 32-degree-of-freedom BFS element is developed, with the same complex Bloch phase applied to displacement, first derivatives, and the mixed derivative.
  \item Across the $(\theta,\mathrm{AR})$ design space, we evaluate directional dispersion, gap measures, and group velocity. The energy-flux relation follows from local balance; the analysis also separates directional stop bands from complete gaps and measures steering amplitude and lobe orientation.
\end{itemize}

Section~\ref{sec:continuum} defines the anisotropic strain-gradient continuum and micro-inertia. Section~\ref{sec:bloch} sets out the Bloch--Floquet formulation, followed by the BFS discretization in Section~\ref{sec:fem}. Section~\ref{sec:verification} reports the analytical, numerical, and mesh checks. Section~\ref{sec:results} presents the homogeneous dispersion, directional gaps, design maps, steering, and the separate heterogeneous Case~C result; Section~\ref{sec:microinertia} examines the high-wavenumber response. Sections~\ref{sec:discussion} and~\ref{sec:conclusions} interpret the findings and state their limitations.
